\section{Sampling the actual source circuit}
\label{sec:peps_actual_source_sampling}

We complete the Gaussian sampling estimate for an actual source-only circuit in
\cite[proof of Theorem 5.2]{OpenAI2026PolynomialPEPS}.
% Source: 04-compression.tex, lines 342--558.
The source occurrences, local branch labels, physical registers,
and discarded registers are those of the original circuit. The sampled physical
quantity is an operator; positivity is not asserted. The tensor-network
representation of that operator is a separate statement.

\begin{theorem}[Finite memories of the discarded registers]
    \label{thm:peps_finite_register_memories}
    \lean{TNLean.PEPS.PairEffect.Layout.finiteDimensional_mem_of_registers,TNLean.PEPS.PairEffect.Layout.finiteDimensional_mem_restrict_mapOwner_mapOwner}
    \leanok
    Let $R$ be a finite ordered list of registers, each with a finite-dimensional
    complex Hilbert space. Its tensor-product memory is finite dimensional.
    The same is true after relabelling owners twice and selecting either side
    of any division of the resulting owners. Zero-dimensional register spaces
    are allowed.
\end{theorem}
\begin{proof}
    \leanok
    Induct on the register list, using finite dimensionality of tensor products.
    Relabelling changes no register space, and selection retains a sublist.
\end{proof}

\begin{definition}[Orthonormal coordinates of selected discarded registers]
    \label{def:peps_restricted_discard_basis}
    \lean{TNLean.PEPS.PairEffect.Layout.restrictedDiscardBasis}
    \leanok
    Under the assumptions of Theorem~\ref{thm:peps_finite_register_memories},
    choose an orthonormal basis of the selected memory indexed by the integers
    from zero to one less than its dimension. The registerwise finite-dimensional
    assumption supplies this basis; no dimension bound is required.
\end{definition}

\begin{definition}[Exterior physical coordinates]
    \label{def:peps_exterior_family_physical_basis}
    \lean{TNLean.PEPS.PairEffect.exteriorFamilyPhysicalBasis}
    \leanok
    Let $P$ be the finite set of parties, listed without repetition, and
    let party $p$ have physical dimension $d_p$. For $A\subseteq P$, the
    tensor product of the original coordinate bases outside $A$ is indexed
    by $\prod_{p\notin A}\{0,\ldots,d_p-1\}$. Transport this basis through
    the canonical owner identifications to the actual exterior physical
    memory. This is an orthonormal basis of that memory.
\end{definition}

\begin{theorem}[Exact exterior physical dimension]
    \label{thm:peps_exterior_family_physical_dimension}
    \lean{TNLean.PEPS.PairEffect.finrank_exteriorFamilyPhysicalLayout}
    \leanok
    The dimension of the actual exterior physical memory after the two
    owner identifications is $\prod_{p\notin A}d_p$.
\end{theorem}
\begin{proof}
    \leanok
    \uses{def:peps_exterior_family_physical_basis}
    Count the coordinate indices of the exterior orthonormal basis.
\end{proof}

\begin{theorem}[Expected norm of an actual partial branch]
    \label{thm:peps_actual_gaussian_branch}
    \lean{TNLean.PEPS.PairEffect.SourceCircuit.integrable_actualGaussianBranch_and_integral_le}
    \leanok
    Let $W$ be an allowed circuit with output memory consisting of physical
    registers followed by discarded registers. Fix a set $S$ of original
    source occurrences and let $A(S)$ be their endpoint parties. Let $p_0$
    be an allowed scalar-input composition containing no sources, and fix
    two actual partial branches $\xi,\zeta$ for $A(S)$.

    At every original occurrence and local branch label, let the two
    endpoint matrices have orthonormal columns, and let the source weights
    be nonnegative with sum one. Choose orthonormal bases for the affected
    and exterior physical memories, indexed by $X,Y$, and for the affected
    and exterior discarded memories, indexed by $D,E$. These four bases,
    transported through the two owner relabellings and canonical register
    partitions, determine bases of the original physical and discarded
    memories.

    For $k>0$, let $M_{\xi,\zeta}(\omega)$ be the actual Gaussian partial
    branch operator on the input $p_0(1)$, traced over the discarded memory
    in these bases. Then its trace norm is integrable for the original
    source Gaussian law and
    \begin{align}
        \mathbb E\,\|M_{\xi,\zeta}\|_1
         & \leq |X|^2 k^{-|S|/2}.
        \label{eq:peps_actual_branch_bound}
    \end{align}
\end{theorem}
\begin{proof}
    \leanok
    \uses{thm:peps_actual_corrected_schmidt_output,thm:peps_gaussian_separated_physical_bound,thm:peps_local_schmidt_output_coordinates}
    Construct the actual local-word Schmidt output representation
    independently for $\xi$ and $\zeta$. The canonical output-coordinate
    identities identify each original prepared output vector with the
    separated physical coordinates of those local words. Reindex the two
    corrected-coordinate sums by the original occurrences in $S$.
    Thus the original partial branch operator is exactly the separated
    physical operator with the original Gaussian coefficient products.
    Its established integrability and Gaussian estimate give
    (\ref{eq:peps_actual_branch_bound}).
\end{proof}

\begin{theorem}[Expected norm of the actual corrected term]
    \label{thm:peps_actual_corrected_source_bound}
    \lean{TNLean.PEPS.PairEffect.SourceCircuit.integral_actualCorrectedSourceTerm_le}
    \leanok
    Let $W$ be an allowed circuit, with output registers consisting of
    physical registers followed by discarded registers. Let $p_0$ be an
    allowed scalar-input composition containing no sources, and use
    $|p_0(1)\rangle\langle p_0(1)|$ as the input. Fix a set $S$ of original
    source occurrences and actual Schmidt decompositions of every source
    vector, with nonnegative weights of sum one and orthonormal endpoint
    columns.

    Let $A(S)$ be the endpoint parties of $S$. Choose orthonormal bases of
    the affected and exterior physical memories and of the affected and
    exterior discarded memories after the canonical owner identifications.
    Write $X$ for the affected physical basis index. For $k>0$, the actual
    physical corrected term $F_S$, in any original physical and discarded
    orthonormal bases, satisfies
    \begin{align}
        \mathbb E\,\|F_S\|_1
         & \leq
        \left(\sum_{\xi,\zeta}
        |c_{\xi}\overline{c_{\zeta}}|\right)
        |X|^2 k^{-|S|/2}.
        \label{eq:peps_actual_corrected_bound}
    \end{align}
    Here $\xi,\zeta$ are the actual partial choices at gate occurrences
    touching $A(S)$, and $c_\xi,c_\zeta$ are their original coefficients.
\end{theorem}
\begin{proof}
    \leanok
    \uses{thm:peps_actual_gaussian_branch,thm:peps_physical_corrected_basis_invariance,thm:peps_gaussian_branch_density_bound}
    Express the actual corrected term in the bases adapted to $A(S)$.
    Trace norm is unchanged by these physical and discarded basis changes.
    Expand the term into its actual partial branches, apply the expectation
    triangle inequality with the factors
    $|c_\xi\overline{c_\zeta}|$, and use
    Theorem~\ref{thm:peps_actual_gaussian_branch} for every pair. The
    integrability needed to pass through the finite sums follows from the
    original Gaussian coefficient products. Factor the common bound out
    of the sum.
\end{proof}

\begin{theorem}[Corrected-term bound from whole-lifetime participation]
    \label{thm:peps_actual_source_lifetime_bound}
    \lean{TNLean.PEPS.PairEffect.SourceCircuit.integral_actualCorrectedSourceTerm_le_lifetime}
    \leanok
    Let the finite set of parties be listed without repetition, with one
    physical register of dimension $d_p$ at each party. Assume every
    discarded register is finite dimensional; no upper bound on these
    dimensions is imposed. Let $W$ be an allowed circuit and let $p_0$
    be an allowed source-free scalar-input preparation.

    Suppose each party participates in at most $b$ gate occurrences over
    the whole circuit. At every gate occurrence, suppose the sum of the
    absolute values of its branch coefficients is at most $B$, and suppose
    $d_p\leq D$ at every party, where $B,D\geq1$. Fix actual normalized
    Schmidt decompositions of all source vectors. Then for every set $S$
    of original source occurrences and every $k>0$, in arbitrary original
    input, physical output and discarded orthonormal coordinates,
    \begin{align}
        \mathbb E\,\|F_S\|_1
         & \leq
        \bigl(B^{4b}D^4\bigr)^{|S|} k^{-|S|/2}.
        \label{eq:peps_actual_lifetime_bound}
    \end{align}
    The input is $|p_0(1)\rangle\langle p_0(1)|$ and the expectation is
    with respect to the original source Gaussian law.
\end{theorem}
\begin{proof}
    \leanok
    \uses{thm:peps_actual_corrected_source_bound,thm:peps_actual_source_lifetime_cost,thm:peps_affected_original_dimensions,def:peps_exterior_family_physical_basis}
    Use the original affected and complementary physical tensor bases.
    Finite dimensionality of each discarded register supplies the two
    regional discarded bases. Their dimensions do not enter
    (\ref{eq:peps_actual_corrected_bound}). The squared affected physical
    dimension is $\prod_{p\in A(S)}d_p^2$. The whole-lifetime incidence
    estimate bounds the product of this quantity with the actual ket--bra
    coefficient sum by $\bigl(B^{4b}D^4\bigr)^{|S|}$. Multiply by the
    nonnegative factor $k^{-|S|/2}$.
\end{proof}

\begin{theorem}[Choosing one sample from corrected-term estimates]
    \label{thm:peps_source_sample_choice_conditional}
    \lean{TNLean.PEPS.PairEffect.SourceCircuit.exists_sampledSourceMatrix_error_le_of_integral_bound}
    \leanok
    Fix a source-only circuit with finite input, physical-output and
    discarded-output orthonormal coordinates, and an arbitrary input matrix.
    Fix endpoint data reproducing every original source vector, with
    nonnegative source weights. Let $N$ be the number of original source
    occurrences, let $Q\geq0$, and let $0<\varepsilon\leq1$. Put
    \[
        k=\left\lceil\left(\frac{8\max(1,N)Q}{\varepsilon}\right)^2\right\rceil+1.
    \]
    Suppose that for every nonempty set $S$ of original source occurrences,
    \[
        \mathbb E\,\|F_S\|_1\leq Q^{|S|}k^{-|S|/2}.
    \]
    Then there is one actual global Gaussian sample for which the sampled
    physical operator differs from the original physical operator in trace norm
    by at most $\varepsilon/4$.
\end{theorem}
\begin{proof}
    \leanok
    \uses{thm:peps_source_sample_count_bound,thm:peps_sampled_error_integral_bound}
    Integrability follows from the finite Gaussian expressions for the actual
    source corrections. The exact nonempty-subset expansion bounds the mean
    error by the sum of the displayed expectations. Choose a sample with error
    at most that sum, and apply the scalar sample-count estimate. This implication
    assumes the corrected-term estimates; the preceding circuit estimate
    supplies them in the application below.
\end{proof}

\begin{theorem}[A single sample for the actual source circuit]
    \label{thm:peps_actual_source_sampling}
    \lean{TNLean.PEPS.PairEffect.SourceCircuit.exists_sourceSample_lifetime_error_le}
    \leanok
    Let $W$ be an allowed circuit on a finite set of parties, listed without
    repetition, with physical output $\bigotimes_p\mathbb C^{d_p}$ followed
    by discarded registers. Every discarded register is finite dimensional,
    with no bound imposed on its dimension. Let $p_0$ be an allowed
    scalar-input composition containing no sources, and use
    $|p_0(1)\rangle\langle p_0(1)|$ as input.

    Suppose every party participates in at most $b$ gate occurrences over
    the whole circuit, every gate has sum of absolute branch coefficients
    at most $B$, and $d_p\leq D$, where $B,D\geq1$. Let $N$ be the number
    of original source occurrences. For $0<\varepsilon\leq1$, put
    \begin{align}
        Q & = B^{4b}D^4,
          & k            & = \left\lceil
        \left(\frac{8\max(1,N)Q}{\varepsilon}\right)^2
        \right\rceil+1.
        \label{eq:peps_actual_sample_count}
    \end{align}
    There are normalized Schmidt data for every original source occurrence
    and local branch label, and one global Gaussian sample $\omega$ using
    $k$ samples for each original occurrence and ordered local label pair,
    such that
    \begin{align}
        \|\widehat\rho(\omega)-\rho\|_1 & \leq \varepsilon/4.
        \label{eq:peps_actual_sample_error}
    \end{align}
    Here $\rho$ is the original physical output operator, and
    $\widehat\rho(\omega)$ is obtained by replacing each mixed source
    operator by its averaged Gaussian operator and tracing the original
    discarded registers. The conclusion holds in arbitrary original input,
    physical and discarded orthonormal coordinates.
\end{theorem}
\begin{proof}
    \leanok
    \uses{thm:peps_actual_source_lifetime_bound,thm:peps_local_schmidt_sources,thm:peps_source_sample_choice_conditional}
    Choose the Schmidt data of all original source occurrences and local
    labels once. For every nonempty corrected set $S$, the actual lifetime
    estimate gives $\mathbb E\|F_S\|_1\leq Q^{|S|}k^{-|S|/2}$.
    The subset expansion and integrability of the Gaussian corrections
    give a single sample whose error is at most the sum of these expectations.
    The choice~(\ref{eq:peps_actual_sample_count}) bounds that sum by
    $\varepsilon/4$. The same sample supplies all source replacements in
    (\ref{eq:peps_actual_sample_error}).
\end{proof}

\begin{theorem}[Explicit power bounds for the sample count]
    \label{thm:peps_source_sampling_polynomial}
    \lean{TNLean.PEPS.Approximation.sourceSamplingCount_lt_of_power_bounds,TNLean.PEPS.Approximation.sourceSamplingCount_lt_inv_pow,TNLean.PEPS.Approximation.sourceSamplingCount_lt_of_gate_power_bounds}
    \leanok
    Let $L\geq1$, $Q\geq0$, $\varepsilon>0$, and suppose
    \[
        N\leq C_NL^n,\qquad Q\leq C_QL^q,\qquad
        \varepsilon^{-1}\leq C_\varepsilon L^p,
    \]
    where $n,q,p$ are nonnegative integers, $C_N\geq1$, and $C_Q\geq0$.
    For the sample count above,
    \[
        k<\bigl(64C_N^2C_Q^2C_\varepsilon^2+2\bigr)
        L^{2(n+q+p)}.
    \]
    At $\varepsilon=L^{-p}$ one may take $C_\varepsilon=1$.
    In particular, if $B,D\geq0$, $B\leq C_BL^m$, $D\leq C_DL^r$,
    and $C_B\geq0$, then $Q=B^{4b}D^4$ gives
    \[
        k<\bigl[64C_N^2(C_B^{4b}C_D^4)^2+2\bigr]
        L^{2(n+4bm+4r+p)}.
    \]
    These inequalities include the case of no source occurrences.
\end{theorem}
\begin{proof}
    \leanok
    \uses{thm:peps_source_sample_count_bound}
    The source bound also bounds $\max(1,N)$ because $C_NL^n\geq1$.
    Substitute the three power bounds into the defining square and use
    $k<(8\max(1,N)Q/\varepsilon)^2+2$. Since the resulting power of $L$
    is at least one, the additive constant is absorbed by the displayed
    coefficient. For $Q=B^{4b}D^4$, multiply the two power estimates.
\end{proof}

\begin{theorem}[Polynomial sampling for arbitrary real accuracy exponents]
    \label{thm:peps_actual_source_polynomial}
    \lean{TNLean.PEPS.PairEffect.SourceCircuit.exists_sourceSample_polynomial_error_le}
    \leanok
    Under the circuit and input assumptions of
    Theorem~\ref{thm:peps_actual_source_sampling}, let $L\geq1$ and suppose
    \begin{align}
        N & \leq C_N L^n,
          & B             & \leq C_B L^m,
          & D             & \leq C_D L^r,
        \label{eq:peps_actual_resource_powers}
    \end{align}
    where $n,m,r$ are nonnegative integers, $C_N\geq1$, and $C_B\geq0$.
    Fix any real accuracy exponent $p$, and let
    $q=\max(0,\lceil p\rceil)$. Choose exactly the sample count
    (\ref{eq:peps_actual_sample_count}) at accuracy $L^{-q}$.
    This positive integer satisfies
    \begin{align}
        k & < \left[64C_N^2\bigl(C_B^{4b}C_D^4\bigr)^2+2\right]
        L^{2(n+4bm+4r+q)}.
        \label{eq:peps_actual_polynomial_count}
    \end{align}
    There are actual global Schmidt data and a sample at this same $k$
    for which $\|\widehat\rho-\rho\|_1\leq L^{-p}/4$.
    In particular this covers every real $p>0$ prescribed in the source.
\end{theorem}
\begin{proof}
    \leanok
    \uses{thm:peps_actual_source_sampling,thm:peps_source_sampling_polynomial}
    Since $L\geq1$ and $q\geq p$, the accuracy $L^{-q}$ belongs to $(0,1]$
    and satisfies $L^{-q}\leq L^{-p}$. Apply the actual sampling theorem
    at $L^{-q}$, retaining its Schmidt data and its sample. The scalar
    sample-count estimate applied to~(\ref{eq:peps_actual_resource_powers})
    gives~(\ref{eq:peps_actual_polynomial_count}) for that exact integer.
\end{proof}
\begin{theorem}[Polynomial stack length at the circuit error budget]
    \label{thm:peps_effect_stack_polynomial}
    \lean{TNLean.PEPS.PairEffect.stackLength_sourceGateBudget_eq}
    \lean{TNLean.PEPS.PairEffect.stackLength_sourceGateBudget_lt_of_power_bounds}
    \leanok
    Let $r,N$ be nonnegative integers, let $S$ be real, and let $\varepsilon>0$.
    Use the per-gate error budget
    $\delta=\varepsilon/(8\max(1,N))$. The chosen stack length is exactly
    \begin{align}
        m & = \left\lceil(rS/\delta)^2\right\rceil+1
        = \left\lceil
        \bigl(8\max(1,N)rS/\varepsilon\bigr)^2
        \right\rceil+1.
        \label{eq:peps_effect_stack_count}
    \end{align}
    Suppose $L\geq1$, $N\leq C_NL^n$ and $0\leq S\leq C_SL^s$, where
    $C_N\geq1$, $C_S\geq0$, and $n,s,p$ are nonnegative integers. At
    $\varepsilon=L^{-p}$, put $H=64C_N^2(rC_S)^2+2$. Then
    $m<H L^{2(n+s+p)}$.
\end{theorem}
\begin{proof}
    \leanok
    \uses{thm:peps_source_sampling_polynomial}
    Substitute the per-gate budget into the stack-length formula. It is
    the same ceiling expression as the source sample count, with cost
    $rS$. Apply the inverse-power sample-count estimate to
    $rS\leq rC_SL^s$.
\end{proof}

\begin{theorem}[Polynomial number of actual replacement monomials]
    \label{thm:peps_effect_replacement_polynomial}
    \lean{TNLean.PEPS.PairEffect.replacementMonomialBound_le_of_power_bounds}
    \lean{TNLean.PEPS.PairEffect.OriginalCircuit.card_branchLabels_replacement_le_power}
    \leanok
    Under the preceding power bounds, let $K$ be a nonnegative integer and
    suppose $K\leq C_KL^a$, with
    $C_K\geq0$ and $a$ a nonnegative integer. Then
    \begin{align}
        K m^r & \leq C_K H^r L^{a+2r(n+s+p)}.
        \label{eq:peps_effect_monomial_count}
    \end{align}
    If every original nonprivate gate has at most $K$ monomials, at most
    $r$ pair effects in each monomial, and absolute coefficient sum at
    most $S$, every gate occurrence of the actual source-only replacement
    has at most the number in~(\ref{eq:peps_effect_monomial_count}) of
    branch labels. Original occurrences remain distinct. The bound also
    includes empty gate expansions and $r=0$.
\end{theorem}
\begin{proof}
    \leanok
    \uses{thm:peps_effect_stack_polynomial}
    Raise the nonnegative stack bound to the fixed integer power $r$ and
    multiply by the bound for $K$. The replacement construction has at
    most $Km^r$ monomials at each original gate occurrence; inserted
    private operations and register permutations introduce no expanded
    gates. Apply this bound at the chosen occurrence.
\end{proof}
